\clearpage\begingroup
\NoSourceHead\SourceLeaf{122}
\makeatletter\tagsleft@true\let\veqno\@@leqno\makeatother
\begin{center}
{\Large [DIOPTRIK.]}\par\bigskip
{\large\textls[150]{NACHLASS.}}\par\bigskip
[I.]\par\medskip
NEUER ALGORITHMUS.\par\medskip
\rule{\linewidth}{.4pt}\par
{\small [Aus Handbuch 21, Bg, Aufsätze, Notizen und Rechnungen zur Mathematik und Astronomie gehörig; Angefangen September 1813, S. 27, 28.]}\par
\rule{\linewidth}{.4pt}
\end{center}
\noindent\hspace*{\parindent}Die Grössen $A, B, C, D$ etc. entstehen aus $a, b, c, d$ etc. so, dass
\[
A=a,\quad B=bA-1,\quad C=cB-A,\quad D=dC-B\quad\text{etc.}
\]
\noindent Wir bezeichnen dann
\[
A\text{ mit }(a,a),\quad B\text{ mit }(a,b),\quad
C\text{ mit }(a,c),\quad D\text{ mit }(a,d)\quad\text{etc.},
\]
\noindent wo also bei $a, b, c, d$ etc. die Ordnung als bestimmt betrachtet wird.

Dieser Algorithmus bietet die interessante Eigenschaft, dass man\textsuperscript{[*)]}\begingroup\renewcommand{\thefootnote}{}\footnotetext{[*) In der Handschrift steht hier noch das Wort »statt«.]}\endgroup\ \textls[120]{vier} oder mehr aufeinander folgende Elemente immer durch \textls[120]{drei} ersetzen kann.

Man hat dann folgende schöne Lehrsätze (Ordn[un]g $a,b,c,\, .\, .\, .\ g,h,i,\, .\, .\, .\ p,q,r,\, .\, .\, .\ z$)
\begin{equation}
(a,p)\mathbin{.}(i,z)-(i,p)\mathbin{.}(a,z)=(a,g)\mathbin{.}(r,z),,
\tag*{I.}
\end{equation}
\begin{equation}
\begin{aligned}
\partial(a,z)={}&(b,z)\partial a+(a,a)(c,z)\partial b\mathbin{+}{}\\
& (a,b)(d,z)\partial c+(a,c)(e,z)\partial d+\text{etc.}
\end{aligned}
\tag*{II.}
\end{equation}

\clearpage\SourceLeaf{123}\HistoricalHead{118}{L}{PHYSIK.\quad NACHLASS.}
Diese Formeln werden noch eleganter, wenn man die Bezeichnung\textsuperscript{[*)]}\begingroup\renewcommand{\thefootnote}{}\footnotetext{[*) In der Handschrift steht hier noch das Wort »anders«.]}\endgroup\ modificirt und
\[
\begin{gathered}
(a,a)=0,\quad (a,b)=1,\quad (a,c)=b,\\
(a,d)=c(a,c)-(a,b),\quad (a,e)=d(a,d)-(a,c)\quad\text{etc.}
\end{gathered}
\]
\noindent setzt. Alsdann hat man nemlich:
\begin{equation}
(a,p)\mathbin{.}(i,z)-(i,p)\mathbin{.}(a,z)=(a,i)\mathbin{.}(p,z),
\tag*{I[a].}
\end{equation}
\begin{equation}
\begin{aligned}
\partial(a,z)={}&(a,b)(b,z)\partial b+(a,c)(c,z)\partial c\mathbin{+}{}\\
& (a,d)(d,z)\partial d+\text{etc.}
\end{aligned}
\tag*{II[a].}
\end{equation}

Auf Einer Axe seyn $\mu+1$ Gläser geordnet, deren Brennweiten $\frac{1}{f},\frac{1}{f'},\frac{1}{f''},\frac{1}{f'''}$ etc. $\frac{1}{f^\mu}$, die Distanzen von einander $h,h',h'',\, .\, .\, .\ h^{\mu-1}$ sind. Man mache, indem obigen Elementen $a,b,c$ etc. jetzt die Grössen $f,h,f',h',f'',h''$ etc. entsprechen:
\[
\begin{array}{cc}
[\mathrm{III.}]\left\{\begin{aligned}
\alpha&=1,\\ \alpha'&=(f,h),\\ \alpha''&=(f,h'),\\ \alpha'''&=(f,h''),
\end{aligned}\right.
&[\mathrm{IV.}]\left\{\begin{aligned}
\beta&=f,\\ \beta'&=(f,f'),\\ \beta''&=(f,f''),\\ \beta'''&=(f,f'''),
\end{aligned}\right.\\[1em]
{}[\mathrm{V.}]\left\{\begin{aligned}
\gamma&=0,\\ \gamma'&=h,\\ \gamma''&=(h,h'),\\ \gamma'''&=(h,h''),
\end{aligned}\right.
&[\mathrm{VI.}]\left\{\begin{aligned}
\delta&=1,\\ \delta'&=(h,f'),\\ \delta''&=(h,f''),\\ \delta'''&=(h,f''')
\end{aligned}\right.
\end{array}
\]
\begin{center}etc.\end{center}
\noindent\hspace*{\parindent}Die Bedeutung dieser Grössen versinnlicht folgende Figur:
\begin{center}
\begin{tikzpicture}[x=.20pt,y=-.38pt,line width=.35pt,every node/.style={font=\scriptsize,inner sep=.5pt}]
\path[use as bounding box] (30,10) rectangle (2055,475);
\draw (91,253)--(1995,262);
\draw (91,253)--(425,165)--(847,364)--(1268,141)--(1637,391)--(1961,111);
\draw (97,119)--(97,255) (425,164)--(425,337) (847,108)--(847,365) (1269,142)--(1269,368) (1637,139)--(1637,391);
\node at (100,50) {Objectiv};
\node[fill=white] at (286,260) {$\frac{\gamma'}{\delta}$};
\node[fill=white] at (503,269) {$\frac{\gamma'}{\delta'}$};
\node[fill=white] at (754,264) {$\frac{\gamma''}{\delta'}$};
\node[fill=white] at (938,254) {$\frac{\gamma''}{\delta''}$};
\node[fill=white] at (1179,266) {$\frac{\gamma'''}{\delta''}$};
\node[fill=white] at (1337,270) {$\frac{\gamma'''}{\delta'''}$};
\node[fill=white] at (1554,260) {$\frac{\gamma^{\mathrm{IV}}}{\delta'''}$};
\node[fill=white] at (1709,260) {$\frac{\gamma^{\mathrm{IV}}}{\delta^{\mathrm{IV}}}$};
\node at (205,225) {$\delta$};
\node at (548,226) {$\delta'$};
\node at (996,287) {$\delta''$};
\node at (1393,226) {$\delta'''$};
\node at (1844,224) {$\delta^{\mathrm{IV}}$};
\node[anchor=west] at (441,218) {$\gamma'$};
\node[anchor=west] at (863,301) {$\gamma''$};
\node[anchor=west] at (1285,201) {$\gamma'''$};
\node[fill=white] at (1650,326) {$\gamma^{\mathrm{IV}}$};
\node at (1040,450) {[Fig. A.]};
\end{tikzpicture}

\begin{tikzpicture}[x=.20pt,y=-.38pt,line width=.35pt,every node/.style={font=\scriptsize,inner sep=.5pt}]
\path[use as bounding box] (25,0) rectangle (2055,410);
\draw (84,191)--(2037,192);
\draw (84,111)--(420,306)--(845,63)--(1265,292)--(1632,20)--(1891,340);
\draw (84,110)--(84,310) (420,65)--(420,305) (844,62)--(844,344) (1264,58)--(1264,293) (1632,20)--(1632,344);
\node[fill=white] at (124,239) {$\frac{\alpha}{\beta}$};
\node[fill=white] at (325,184) {$\frac{\alpha'}{\beta}$};
\node[fill=white] at (489,183) {$\frac{\alpha'}{\beta'}$};
\node[fill=white] at (746,204) {$\frac{\alpha''}{\beta'}$};
\node[fill=white] at (920,205) {$\frac{\alpha''}{\beta''}$};
\node[fill=white] at (1173,180) {$\frac{\alpha'''}{\beta''}$};
\node[fill=white] at (1313,177) {$\frac{\alpha'''}{\beta'''}$};
\node[fill=white] at (1530,205) {$\frac{\alpha^{\mathrm{IV}}}{\beta'''}$};
\node[fill=white] at (1694,209) {$\frac{\alpha^{\mathrm{IV}}}{\beta^{\mathrm{IV}}}$};
\node at (164,165) {$\beta$};
\node at (555,222) {$\beta'$};
\node at (1003,155) {$\beta''$};
\node at (1380,224) {$\beta'''$};
\node at (1833,228) {$\beta^{\mathrm{IV}}$};
\node[anchor=east] at (68,162) {$\alpha$};
\node[anchor=east] at (404,244) {$\alpha'$};
\node[anchor=west] at (860,130) {$\alpha''$};
\node[anchor=east] at (1248,228) {$\alpha'''$};
\node[anchor=east] at (1616,115) {$\alpha^{\mathrm{IV}}$};
\node at (1040,382) {[Fig. B.]};
\end{tikzpicture}
\end{center}


\clearpage\SourceLeaf{124}\HistoricalHead{119}{R}{NEUER ALGORITHMUS.}
Hier wird nun:
\par\smallskip\noindent
\begin{tabular}{@{}p{.29\linewidth}@{\hspace{1em}}p{\dimexpr.71\linewidth-1em\relax}@{}}
\multicolumn{2}{@{}p{\linewidth}@{}}{$\delta^\mu$ die Vergrösserung; bei geradem $\mu$ ist $\delta^\mu$ positiv bei aufrechtem Bilde, negativ bei verkehrtem,}\\[.5em]
$\frac{g}{\beta},\ \frac{g}{\beta'},\ \frac{g}{\beta''}$ etc.&Grösse der einzelnen Bilder, wenn $g$ das Gesichtsfeld,\\[.5em]
$g\gamma',\ g\gamma'',\ g\gamma'''$ etc.&Oeffnungen der Gläser wegen des Gesichtsfeldes,\\[.5em]
$\frac{\omega}{\delta},\ \frac{\omega}{\delta'},\ \frac{\omega}{\delta''}$ etc.&Durchmesser der Stralenkegel,\\[.5em]
$\omega\alpha',\ \omega\alpha'',\ \omega\alpha'''$ etc.&Oeffnungen der Gläser wegen der Helligkeit.
\end{tabular}
\par\medskip
Für weitsichtige Augen muss $\beta^\mu=0$ werden; für kurzsichtige, deren Augenweite $=k$, muss $\beta^\mu\delta^\mu k=-1$ seyn.

Die Breite des farbigen Saumes um das Gesichtsfeld wird
\begin{equation}
\tfrac{1}{2}\delta^\mu g(\alpha^\mu\partial\delta^\mu-\beta^\mu\partial\gamma^\mu)
\tag*{[VII]}
\end{equation}
\noindent und zwar ist der Saum, wenn diese Grösse positiv wird, blau; roth, wenn sie negativ ist
\begin{equation}
\begin{aligned}
={}&\frac{1}{2}\delta^\mu g\frac{\partial n}{n-1}\{\alpha'\gamma'f'+\alpha''\gamma''f''\mathbin{+}{}\\
&\alpha'''\gamma'''f'''+\alpha^{\mathrm{IV}}\gamma^{\mathrm{IV}}f^{\mathrm{IV}}+\text{etc.}+\alpha^\mu\gamma^\mu f^\mu\},
\end{aligned}
\tag*{[VIIa]}
\end{equation}
\begin{equation}
\begin{aligned}
={}&\frac{1}{2}\delta^\mu g\frac{\partial n}{n-1}\{\alpha'(\delta+\delta')+\alpha''(\delta'+\delta'')\mathbin{+}{}\\
&\alpha'''(\delta''+\delta''')+\text{etc.}+\alpha^\mu(\delta^{\mu-1}+\delta^\mu)\},
\end{aligned}
\tag*{[VIIb]}
\end{equation}
\begin{equation}
\begin{aligned}
={}&\frac{1}{2}\delta^\mu g\frac{\partial n}{n-1}\{\gamma'(\beta+\beta')+\gamma''(\beta'+\beta'')\mathbin{+}{}\\
&\gamma'''(\beta''+\beta''')+\text{etc.}+\gamma^\mu(\beta^{\mu-1}+\beta^\mu)\}\text{[*)]}.\!
\end{aligned}
\tag*{[VIIc]}
\end{equation}
\begingroup\renewcommand{\thefootnote}{}\footnotetext{[*) In der Handschrift steht in den drei letzten Formeln $\delta^m$ statt $\delta^\mu$.]}\endgroup
\begin{center}\rule{.25\linewidth}{.4pt}\end{center}
\clearpage\endgroup
